\begin{abstract}
Global floating photovoltaic (FPV) assessments need to distinguish a declared resource scenario from evidence-supported screening opportunity and project deployability. We reconciled 199{,}976 large lakes and mapped reservoirs, applied a common monthly electricity model, and propagated explicit evidence states through a reference resource scenario (L0), an ice screen (L1), public-evidence suitability bounds (L2), and a partial infrastructure-proximity screen (L3). L0 was 16{,}713~TWh~yr$^{-1}$ under 30\% reservoir coverage, 10\% lake coverage, and a 30~km$^2$ footprint cap; L1 retained 10{,}995~TWh~yr$^{-1}$ (65.79\%). Core and conservative L2 intervals retained 4{,}190--4{,}805 and 2{,}914--3{,}404~TWh~yr$^{-1}$, respectively (17.44--28.75\% of L0). Adding bounded road and mapped-power-line proximity produced core and conservative partial-L3 intervals of 1{,}595--3{,}508 and 1{,}101--2{,}494~TWh~yr$^{-1}$ (6.58--20.99\%). The specific-yield screen correlated with an independent published lake model ($r=0.896$), while coverage, ice evidence, and infrastructure mapping remained material uncertainty sources. Lower interval endpoints support higher-confidence global screening; upper endpoints identify priorities for data acquisition. Complete suitability under the published 1991--2020 criterion and complete deployable potential remain unavailable.
\end{abstract}
